% GENERATED FILE. Edit canonical lessons, then run npm run generate:books.
% canonical-source-hash: fnv1a64:7069c03a04ef798c
% canonical-lessons: TE-C396-ayinappatiki, TE-C396-marovaipu, TE-C396-kaga, TE-C396-viruddham, TE-C396-adevidhanga

\chapter{On the Other Hand}
\label{ch:te-on-the-other-hand}
\hlchaptermodality{396}

\begin{chapteropening}
\textbf{By the end of this chapter:} \emph{I can contrast two sides of an argument.}

Complete the last lesson of chapter 396: I can contrast two sides of an argument.
\end{chapteropening}

\section[\texorpdfstring{\te{అయినప్పటికీ} (ayinappaṭikī)}{అయినప్పటికీ (ayinappaṭikī)}]{\te{అయినప్పటికీ} (ayinappaṭikī) — nevertheless; even so}

\label{lesson:TE-C396-ayinappatiki}



\begin{quote}
Before the new one: say the Telugu for that is to say, then the Telugu for a footing.
\end{quote}

\subsection*{You'll want to know: \te{అయినప్పటికీ}}
\textbf{\te{అయినప్పటికీ}} — \emph{ayinappaṭikī} — \textquotedblleft{}nevertheless; even so\textquotedblright{}.

\begin{quote}
\te{అతను} \te{అలసిపోయాడు}, \te{అయినప్పటికీ} \te{పని} \te{పూర్తి} \te{చేశాడు}.
\end{quote}

\begin{quote}
\emph{atanu alasipōyāḍu, ayinappaṭikī pani pūrti cēśāḍu.}
\end{quote}

\begin{quote}
He was tired; even so, he finished the work.
\end{quote}

\te{అయిన}, having been, + \te{అప్పటికీ}, even then. Stronger than \te{అయినా}, which you know.

\begin{cousinweb}
\textbf{\te{అయినప్పటికీ}} (\emph{ayinappaṭikī}) is native Dravidian: \emph{ayina}, a form of \emph{avu}, \textquotedblleft{}to be, to become\textquotedblright{} (DEDR 333), and \emph{appaṭikī}, \textquotedblleft{}even at that time\textquotedblright{}.
\end{cousinweb}

\begin{grammarlens}[title={What you've built}]
One more word for giving reasons and opinions.
\end{grammarlens}

\subsection*{Your turn}
\begin{itemize}
\raggedright
  \item \emph{Say it:} \emph{ayinappaṭikī}
  \item \emph{Say it:} \emph{ayinappaṭikī}, once more
  \item \emph{Say it:} \emph{atanu alasipōyāḍu, ayinappaṭikī pani pūrti cēśāḍu.} — He was tired; even so, he finished the work.
  \item \emph{Recall:} say the Telugu for that is to say, then the Telugu for a footing, then say \emph{ayinappaṭikī} again
\end{itemize}

\begin{tcolorbox}[breakable,colback=teal!4,colframe=teal!35!black,title={Before you move on}]
What does \te{అయినప్పటికీ} mean? (\textquotedblleft{}Nevertheless; even so\textquotedblright{}.) Say it once more.
\end{tcolorbox}
\section[\texorpdfstring{\te{మరోవైపు} (marōvaipu)}{మరోవైపు (marōvaipu)}]{\te{మరోవైపు} (marōvaipu) — on the other hand}

\label{lesson:TE-C396-marovaipu}



\begin{quote}
Before the new one: say the Telugu for a footing, then the Telugu for nevertheless.
\end{quote}

\subsection*{You'll want to know: \te{మరోవైపు}}
\textbf{\te{మరోవైపు}} — \emph{marōvaipu} — \textquotedblleft{}on the other hand\textquotedblright{}.

\begin{quote}
\te{ఒక} \te{వైపు} \te{డబ్బు} \te{కావాలి}, \te{మరోవైపు} \te{సమయం} \te{లేదు}.
\end{quote}

\begin{quote}
\emph{oka vaipu ḍabbu kāvāli, marōvaipu samayaṁ lēdu.}
\end{quote}

\begin{quote}
On one side we need money; on the other hand there is no time.
\end{quote}

\te{ఒక} \te{వైపు} ... \te{మరోవైపు} ..., on one hand ... on the other.

\begin{cousinweb}
\textbf{\te{మరోవైపు}} (\emph{marōvaipu}) is a Telugu phrase of \emph{marō}, \textquotedblleft{}another\textquotedblright{}, and \emph{vaipu}, \textquotedblleft{}a side, a direction\textquotedblright{}; the origin of \emph{vaipu} is uncertain.
\end{cousinweb}

\begin{grammarlens}[title={What you've built}]
One more word for giving reasons and opinions.
\end{grammarlens}

\subsection*{Your turn}
\begin{itemize}
\raggedright
  \item \emph{Say it:} \emph{marōvaipu}
  \item \emph{Say it:} \emph{marōvaipu}, once more
  \item \emph{Say it:} \emph{oka vaipu ḍabbu kāvāli, marōvaipu samayaṁ lēdu.} — On one side we need money; on the other hand there is no time.
  \item \emph{Recall:} say the Telugu for a footing, then the Telugu for nevertheless, then say \emph{marōvaipu} again
\end{itemize}

\begin{tcolorbox}[breakable,colback=teal!4,colframe=teal!35!black,title={Before you move on}]
What does \te{మరోవైపు} mean? (\textquotedblleft{}On the other hand\textquotedblright{}.) Say it once more.
\end{tcolorbox}
\section[\texorpdfstring{\te{కాగా} (kāgā)}{కాగా (kāgā)}]{\te{కాగా} (kāgā) — whereas; while}

\label{lesson:TE-C396-kaga}



\begin{quote}
Before the new one: say the Telugu for nevertheless, then the Telugu for on the other hand.
\end{quote}

\subsection*{You'll want to know: \te{కాగా}}
\textbf{\te{కాగా}} — \emph{kāgā} — \textquotedblleft{}whereas; while\textquotedblright{}.

\begin{quote}
\te{అన్న} \te{వైద్యుడు} \te{కాగా}, \te{తమ్ముడు} \te{రైతు}.
\end{quote}

\begin{quote}
\emph{anna vaidyuḍu kāgā, tammuḍu raitu.}
\end{quote}

\begin{quote}
Whereas the elder brother is a doctor, the younger is a farmer.
\end{quote}

It sets two facts side by side; it is a form of \te{అవు}, to be.

\begin{cousinweb}
\textbf{\te{కాగా}} (\emph{kāgā}) is native Dravidian, a form of \emph{avu}, \textquotedblleft{}to be, to become\textquotedblright{} (DEDR 333), whose stem takes the shape \emph{kā-} in some forms.
\end{cousinweb}

\begin{grammarlens}[title={What you've built}]
One more word for giving reasons and opinions.
\end{grammarlens}

\subsection*{Your turn}
\begin{itemize}
\raggedright
  \item \emph{Say it:} \emph{kāgā}
  \item \emph{Say it:} \emph{kāgā}, once more
  \item \emph{Say it:} \emph{anna vaidyuḍu kāgā, tammuḍu raitu.} — Whereas the elder brother is a doctor, the younger is a farmer.
  \item \emph{Recall:} say the Telugu for nevertheless, then the Telugu for on the other hand, then say \emph{kāgā} again
\end{itemize}

\begin{tcolorbox}[breakable,colback=teal!4,colframe=teal!35!black,title={Before you move on}]
What does \te{కాగా} mean? (\textquotedblleft{}Whereas; while\textquotedblright{}.) Say it once more.
\end{tcolorbox}
\section[\texorpdfstring{\te{విరుద్ధం} (viruddhaṁ)}{విరుద్ధం (viruddhaṁ)}]{\te{విరుద్ధం} (viruddhaṁ) — contrary; opposed}

\label{lesson:TE-C396-viruddham}



\begin{quote}
Before the new one: say the Telugu for on the other hand, then the Telugu for whereas.
\end{quote}

\subsection*{You'll want to know: \te{విరుద్ధం}}
\textbf{\te{విరుద్ధం}} — \emph{viruddhaṁ} — \textquotedblleft{}contrary; opposed\textquotedblright{}.

\begin{quote}
\te{అది} \te{నియమాలకు} \te{విరుద్ధం}.
\end{quote}

\begin{quote}
\emph{adi niyamālaku viruddhaṁ.}
\end{quote}

\begin{quote}
That is against the rules.
\end{quote}

The thing opposed takes -\te{కి}: \te{నియమానికి} \te{విరుద్ధం}.

\begin{cousinweb}
\textbf{\te{విరుద్ధం}} (\emph{viruddhaṁ}) is borrowed from Sanskrit \emph{viruddha}, \textquotedblleft{}opposed, hindered\textquotedblright{} (CDIAL 11864).
\end{cousinweb}

\begin{grammarlens}[title={What you've built}]
One more word for giving reasons and opinions.
\end{grammarlens}

\subsection*{Your turn}
\begin{itemize}
\raggedright
  \item \emph{Say it:} \emph{viruddhaṁ}
  \item \emph{Say it:} \emph{viruddhaṁ}, once more
  \item \emph{Say it:} \emph{adi niyamālaku viruddhaṁ.} — That is against the rules.
  \item \emph{Recall:} say the Telugu for on the other hand, then the Telugu for whereas, then say \emph{viruddhaṁ} again
\end{itemize}

\begin{tcolorbox}[breakable,colback=teal!4,colframe=teal!35!black,title={Before you move on}]
What does \te{విరుద్ధం} mean? (\textquotedblleft{}Contrary; opposed\textquotedblright{}.) Say it once more.
\end{tcolorbox}
\section[\texorpdfstring{\te{అదేవిధంగా} (adēvidhaṅgā)}{అదేవిధంగా (adēvidhaṅgā)}]{\te{అదేవిధంగా} (adēvidhaṅgā) — likewise; in the same way}

\label{lesson:TE-C396-adevidhanga}



\begin{quote}
Before the new one: say the Telugu for whereas, then the Telugu for contrary.
\end{quote}

\subsection*{You'll want to know: \te{అదేవిధంగా}}
\textbf{\te{అదేవిధంగా}} — \emph{adēvidhaṅgā} — \textquotedblleft{}likewise; in the same way\textquotedblright{}.

\begin{quote}
\te{అన్న} \te{బాగా} \te{చదివాడు}, \te{అదేవిధంగా} \te{చెల్లి} \te{కూడా}.
\end{quote}

\begin{quote}
\emph{anna bāgā cadivāḍu, adēvidhaṅgā celli kūḍā.}
\end{quote}

\begin{quote}
The elder brother studied hard, and likewise the younger sister.
\end{quote}

\te{అదే}, that very, + \te{విధం}, a way, + -\te{గా}.

\begin{cousinweb}
\textbf{\te{అదేవిధంగా}} (\emph{adēvidhaṅgā}) is built from native \emph{adē}, \textquotedblleft{}that very\textquotedblright{}, on the far-pointing base \emph{a-} (DEDR 1), and \emph{vidhaṁ}, \textquotedblleft{}a manner\textquotedblright{}, from Sanskrit \emph{vidha} (Monier-Williams), with the adverb ending -\te{గా}.
\end{cousinweb}

\begin{grammarlens}[title={What you've built}]
One more word for giving reasons and opinions.
\end{grammarlens}

\subsection*{Your turn}
\begin{itemize}
\raggedright
  \item \emph{Say it:} \emph{adēvidhaṅgā}
  \item \emph{Say it:} \emph{adēvidhaṅgā}, once more
  \item \emph{Say it:} \emph{anna bāgā cadivāḍu, adēvidhaṅgā celli kūḍā.} — The elder brother studied hard, and likewise the younger sister.
  \item \emph{Recall:} say the Telugu for whereas, then the Telugu for contrary, then say \emph{adēvidhaṅgā} again
\end{itemize}

\begin{tcolorbox}[breakable,colback=teal!4,colframe=teal!35!black,title={Before you move on}]
What does \te{అదేవిధంగా} mean? (\textquotedblleft{}Likewise; in the same way\textquotedblright{}.) Say it once more.
\end{tcolorbox}
